\bookchapter{Capitolo 10: Pagamenti mobili e servizi per la vita quotidiana}

\chapterintro{A dire il vero, i miei genitori e la loro generazione all'inizio erano restii verso “pagare con il telefono”. “I soldi sono nel telefono, cosa succede se me lo rubano?” È un timore normale, ma la realtà è che la probabilità di perdere i contanti in tasca è molto più alta di quella di subire un addebito fraudolento con il pagamento da telefono. Il pagamento mobile non è conservare i soldi nel telefono: è solo una chiave, una chiave che apre il tuo conto bancario.}

\booksection{Come si collega una carta bancaria?}

Che tu usi Alipay o WeChat, la prima volta per collegare la carta devi seguire questi passaggi:

\begin{enumerate}
  \item Apri l'app, trova "Il mio" o "Portafoglio". Tocca "Aggiungi carta bancaria".
  \item Inserisci il numero di carta. Oggi gran parte delle app supporta la fotocamera per riconoscere il numero di carta, senza digitarlo cifra per cifra.
  \item La banca invierà un SMS con codice di verifica al numero di telefono che hai lasciato in banca. Inseriscilo e imposta un PIN di pagamento di 6 cifre.
  \item Questo PIN di pagamento, il PIN della carta bancaria e quello di blocco schermo dovrebbero essere tre password diverse, non impostarli uguali per comodità.
\end{enumerate}

Un piccolo trucco: se sei preoccupato per la sicurezza, puoi aprire nell'app della banca un conto elettronico di classe II o III da dedicare ai pagamenti mobili. Metti su questa carta solo i soldi per le piccole spese quotidiane, così anche in caso di problema le perdite sono limitate.

\booksection{Il codice di pagamento funziona come “chi lo vede, incassa”}

Molti non distinguono tra codice di pagamento e codice di ricezione, ed è una differenza importante.

\begin{itemize}
  \item \textbf{Codice di pagamento}: il codice a barre + QR code che mostri alla cassiera perché lo scansioni. Questo codice equivale al tuo "portafoglio elettronico": chi lo scansiona può prelevare denaro dal tuo conto. Quindi\textbf{non inviare MAI uno screenshot a nessuno}. Il codice si aggiorna automaticamente ogni minuto, proprio per evitare usi fraudolenti.
  \item \textbf{Codice di ricezione}: gli altri lo scansionano per inviarti denaro. Questo codice può essere salvato come screenshot e inviato ad altri, è sicuro.
\end{itemize}

Ricorda questa differenza e potrai evitare molte trappole truffaldine. Chiunque ti chieda di “inviare uno screenshot del codice di pagamento” è al 100\% un truffatore.

\booksection{Queste impostazioni di sicurezza, 5 minuti per configurarle e tanti pensieri in meno}

\begin{itemize}
  \item \textbf{Pagamento con impronta/volto}: nelle impostazioni di pagamento di Alipay e WeChat attiva il pagamento biometrico. In questo modo, a ogni pagamento occorre impronta o volto; anche se perdi il telefono, chi lo trova non può spendere i tuoi soldi.
  \item \textbf{Blocco di sicurezza}: Alipay e WeChat hanno una funzione "Blocco di sicurezza". Una volta attivata, per entrare nella pagina "Portafoglio" dell'app occorre una verifica d'identità aggiuntiva. Percorso: Alipay Il mio Impostazioni Account e sicurezza Centro sicurezza Blocco di sicurezza.
  \item \textbf{Controllo del pagamento senza PIN}: molti hanno sottoscritto vari pagamenti senza PIN e se ne sono dimenticati. Vai nelle "Impostazioni di pagamento" di Alipay e WeChat, trova "Pagamento senza PIN/addebito automatico", controlla i servizi sottoscritti e disattiva quelli non più necessari.
  \item \textbf{Gestione dispositivi di accesso}: controlla periodicamente quali dispositivi hanno effettuato l'accesso al tuo account. Se ne vedi uno sconosciuto, scollegalo immediatamente. Percorso: Alipay Impostazioni Account e sicurezza Gestione dispositivi di accesso.
\end{itemize}

\booksection{Autobus e metro: il telefono è la carta di trasporto}

Oggi, prendendo autobus e metro, sono sempre meno le persone che usano carte di trasporto fisiche. Per usare il telefono ai tornelli ci sono due strade:

\begin{itemize}
  \item \textbf{Carta di trasporto virtuale nel wallet del telefono}: attiva direttamente una carta di trasporto virtuale in Apple Wallet, Huawei Wallet, Xiaomi Wallet. Il vantaggio è che non occorre sbloccare il telefono al passaggio, basta appoggiarlo, la risposta è la più rapida, e funziona anche senza rete. In alcune città la carta virtuale offre gli stessi sconti sui trasbordi di quella fisica.
  \item \textbf{Codice di viaggio di Alipay/WeChat}: cerca "Codice di viaggio" nell'app, attiva il codice di autobus o metro della tua città. Il vantaggio è che non occorre ricaricare, l'importo si deduce dal saldo o dalla carta bancaria collegata. Ma serve la rete e prima di passare devi aprire l'app e trovare il codice.
\end{itemize}

Personalmente consiglio la carta di trasporto virtuale nel wallet del telefono. Anche se la prima volta devi ricaricare 10-20 yuan, è davvero molto più comoda, soprattutto quando si prende la metro di fretta, senza doversi fermare a sbloccare il telefono per cercare l'app.

\booksection{Bollette, consegna cibo a domicilio, taxi}

\begin{itemize}
  \item \textbf{Pagamento bollette (acqua, elettricità, gas)}: in "Pagamenti bollette" di Alipay o WeChat collega i numeri utente di acqua, elettricità e gas di casa tua. Ogni mese, quando arriva la bolletta, basta entrare e pagare. I numeri utente si trovano di solito sulle bollette cartacee di casa, o si possono chiedere al servizio clienti.
  \item \textbf{Confronto dei prezzi del delivery}: lo stesso ristorante su Meituan ed Ele.me ha spesso prezzi e offerte diversi. Prima di ordinare controlla su entrambi, a volte si risparmiano diverse unità. Inoltre la funzione "Pin Hao Fan" di Meituan, ordinando con amici, permette di pagare meno.
  \item \textbf{Sicurezza dei taxi}: in Didi o Amap Ride-hailing, ricordati di impostare "Contatto di emergenza" e "Condivisione percorso". Quando prendi un taxi, l'app condivide automaticamente il tuo percorso con il contatto impostato. Inoltre, di notte privilegia il "Taxi" rispetto a "Kuai Che" (auto privata): il taxi è soggetto a regolamentazione del settore, con un livello di sicurezza più alto.
\end{itemize}

\booksection{Gestire dal telefono assicurazione sanitaria, previdenza e adempimenti amministrativi}

Negli ultimi anni i cambiamenti sono stati enormi: molte pratiche che prima richiedevano code agli sportelli oggi si sbrigano in pochi minuti sul telefono.

\begin{itemize}
  \item \textbf{Credenziale elettronica dell'assicurazione sanitaria}: cerca "Credenziale elettronica assicurazione sanitaria" in Alipay o WeChat. Dopo l'attivazione, in ospedale o in farmacia basta usare il telefono, senza portare la carta fisica. La prima volta serve una verifica con riconoscimento facciale, in seguito a ogni utilizzo viene generato un codice dinamico.
  \item \textbf{Consulta previdenza}: scarica l'app "Palmo 12333" o usa il "Centro cittadino" in Alipay per consultare i versamenti previdenziali, il saldo del conto pensione, richiedere sussidi di disoccupazione, ecc.
  \item \textbf{Imposte sul reddito}: la liquidazione annuale delle imposte sul reddito, da marzo a giugno, si può fare dall'app "Imposte sul reddito personale". Ricordati di controllare se ci sono detrazioni speciali dichiarabili (istruzione dei figli, interessi del mutuo, mantenimento genitori anziani, ecc.), tutte cose che ti permettono di pagare meno tasse o di ottenere un rimborso.
\end{itemize}

\begin{tcbwarn}{Ricorda queste regole per difenderti dalle truffe}
\begin{itemize}
  \item Polizia, magistratura, banche e servizio clienti non ti chiederanno mai di trasferire denaro su un "conto sicuro".
  \item Chiunque ti chieda di inviare uno screenshot del codice di pagamento è un truffatore.
  \item Non cliccare su link sconosciuti, soprattutto in SMS come "ETC scaduto" o "carta previdenziale sospesa".
  \item I rimborsi degli acquisti online tornano sul metodo di pagamento originario, non serve fornire il numero della tua carta bancaria.
\end{itemize}

\end{tcbwarn}

\begin{tcbpractice}{Fallo ora: fai un controllo di sicurezza alla tua app di pagamento}
\begin{enumerate}
  \item Apri Alipay e WeChat, conferma che il pagamento con impronta/volto sia attivo.
  \item Trova la funzione "Blocco di sicurezza" e attivala.
  \item Controlla l'elenco "Pagamento senza PIN/addebito automatico" e disattiva i servizi non più necessari.
  \item Visualizza l'elenco "Dispositivi di accesso" e rimuovi quelli non riconosciuti.
  \item Attiva la tua credenziale elettronica dell'assicurazione sanitaria, così la prossima volta in ospedale non dovrai portare la carta fisica.
\end{enumerate}

\end{tcbpractice}
